\chapter{تصمیم‌ناپذیری}

آیا مسئله‌های واقعی و ملموسی وجود دارند که هیچ کامپیوتری نتواند حلشان کند؟ بله. مثلاً \textbf{درستی‌سنجی برنامه}: یک برنامه و یک توصیف دقیق از کاری که باید انجام دهد داریم و می‌خواهیم مطمئن شویم برنامه دقیقاً همان کار را می‌کند. با اینکه هم برنامه و هم توصیف کاملاً دقیق‌اند، در حالت کلی هیچ الگوریتمی نمی‌تواند این را بررسی کند. در این فصل اولین مسئله تصمیم‌ناپذیر را پیدا می‌کنیم و روش اثبات تصمیم‌ناپذیری را یاد می‌گیریم.

\section{مسئله پذیرش برای ماشین تورینگ}

\[\ATM=\{\enc{M,w}\st M\ \text{\rl{یک ماشین تورینگ است که $w$ را می‌پذیرد}}\}.\]

\begin{theorem}\label{thm:atm}
$\ATM$ تصمیم‌ناپذیر است.
\end{theorem}

قبل از اثبات، ببینیم $\ATM$ \textbf{تشخیص‌پذیر} است:
\begin{algo}{$U$}{روی ورودی $\enc{M,w}$ که $M$ ماشین تورینگ و $w$ رشته است}
\begin{enumerate}
\item $M$ را روی ورودی $w$ شبیه‌سازی کن.
\item اگر $M$ به حالت پذیرش رسید \textbf{بپذیر}؛ اگر به حالت رد رسید \textbf{رد کن}.
\end{enumerate}
\end{algo}
اگر $M$ روی $w$ حلقه بزند، $U$ هم روی $\enc{M,w}$ حلقه می‌زند. پس $U$ تصمیم‌گیرنده نیست. $U$ نمی‌تواند تشخیص دهد که $M$ متوقف نخواهد شد؛ به‌زودی می‌بینیم این کار \textbf{غیرممکن} است.

$U$ نمونه‌ای از \term{ماشین تورینگ جهانی}{Universal Turing Machine} است (تورینگ، ۱۹۳۶): ماشینی که می‌تواند هر ماشین تورینگ دیگری را از روی توصیفش شبیه‌سازی کند. این همان ایده کامپیوترهای «برنامه‌پذیر» امروزی است که برنامه را به عنوان داده می‌گیرند.

\section{روش قطری‌سازی}

برای اثبات قضیه \ref{thm:atm} از روش \textbf{قطری‌سازی} استفاده می‌کنیم که کانتور در سال ۱۸۷۳ ابداع کرد. کانتور می‌خواست اندازه مجموعه‌های نامتناهی را مقایسه کند. نمی‌توان آن‌ها را شمرد، اما می‌توان سعی کرد عضوهایشان را با هم \textbf{جفت} کرد (فصل ۱).

یادآوری: تابع $f:A\to B$ \textbf{یک‌به‌یک} است اگر $a\neq b\rightarrow f(a)\neq f(b)$، \textbf{پوشا} است اگر $\forall b\in B\ \exists a\in A:\ f(a)=b$، و \textbf{تناظر} (دوسویی) است اگر هر دو باشد. $A$ و $B$ \textbf{هم‌اندازه} هستند اگر تناظری بین آن‌ها باشد.

\begin{example}
مجموعه اعداد طبیعی $\mathbb{N}$ و مجموعه اعداد طبیعی زوج $\mathcal{E}$ هم‌اندازه‌اند، چون $f(n)=2n$ یک تناظر بین آن‌هاست. این عجیب به نظر می‌رسد، چون $\mathcal{E}$ زیرمجموعه سره $\mathbb{N}$ است! اما برای مجموعه‌های نامتناهی، «جزء» می‌تواند با «کل» هم‌اندازه باشد.
\end{example}

\begin{definition}[شمارا]
مجموعه $A$ \textbf{شماراست} اگر متناهی باشد یا با $\mathbb{N}$ هم‌اندازه باشد.
\end{definition}

\begin{example}[اعداد گویای مثبت]
مجموعه $\mathcal{Q}=\{\tfrac mn\st m,n\in\mathbb{N}\}$ شماراست. همه کسرها را در یک جدول نامتناهی می‌نویسیم: سطر $i$ کسرهایی با صورت $i$ و ستون $j$ کسرهایی با مخرج $j$. سپس جدول را روی \textbf{قطرهای فرعی} پیمایش می‌کنیم و کسرهای تکراری (مثل $\tfrac22=\tfrac11$) را رد می‌کنیم:
\begin{ltrfig}
\begin{tikzpicture}[scale=0.9,font=\small]
  \foreach \i in {1,...,5} { \foreach \j in {1,...,5} { \node (n\i\j) at (\j*1.1,-\i*0.8) {$\frac{\i}{\j}$}; } }
  \draw[->,jzAccent,thick] (n11) -- (n12); \draw[->,jzAccent,thick] (n12) -- (n21); \draw[->,jzAccent,thick] (n21) -- (n31);
  \draw[->,jzAccent,thick] (n31) -- (n22); \draw[->,jzAccent,thick] (n22) -- (n13); \draw[->,jzAccent,thick] (n13) -- (n14);
  \draw[->,jzAccent,thick] (n14) -- (n23); \draw[->,jzAccent,thick] (n23) -- (n32); \draw[->,jzAccent,thick] (n32) -- (n41);
  \node at (6.6,-2.4) {$\cdots$}; \node at (3.3,-4.6) {$\vdots$};
\end{tikzpicture}
\end{ltrfig}
هر کسر بالاخره روی یکی از قطرها قرار می‌گیرد و هر قطر متناهی است، پس همه کسرها در فهرست می‌آیند: $\tfrac11,\tfrac12,\tfrac21,\tfrac31,\tfrac13,\tfrac14,\tfrac23,\tfrac32,\tfrac41,\dots$
\end{example}

پس آیا همه مجموعه‌های نامتناهی هم‌اندازه‌اند؟ نه. بعضی مجموعه‌ها آن‌قدر بزرگ‌اند که با $\mathbb{N}$ تناظری ندارند. به آن‌ها \textbf{ناشمارا} می‌گوییم.

\begin{theorem}\label{thm:reals}
مجموعه اعداد حقیقی $\mathcal{R}$ ناشماراست.
\end{theorem}
\begin{proof}
با برهان خلف فرض کنید تناظری مثل $f:\mathbb{N}\to\mathcal{R}$ وجود دارد. پس می‌توان همه اعداد حقیقی را در یک فهرست نوشت: $f(1),f(2),f(3),\dots$. عددی مثل $x\in(0,1)$ می‌سازیم که در این فهرست \textbf{نیست}: رقم $i$ام بعد از ممیز $x$ را طوری انتخاب می‌کنیم که با رقم $i$ام بعد از ممیز $f(i)$ فرق داشته باشد. مثلاً:
\begin{ltrfig}
\begin{tabular}{c|l}
$n$ & $f(n)$\\ \hline
$1$ & $3.\underline{1}4159\ldots$\\
$2$ & $55.5\underline{5}555\ldots$\\
$3$ & $0.12\underline{3}45\ldots$\\
$4$ & $0.500\underline{0}0\ldots$\\
$\vdots$ & $\vdots$
\end{tabular}
\qquad\qquad $x=0.4641\ldots$
\end{ltrfig}
(رقم اول $x$ با ۱، رقم دوم با ۵، رقم سوم با ۳ و رقم چهارم با ۰ فرق دارد.) برای جلوگیری از مشکل نمایش‌های دوگانه مثل $0.1999\ldots=0.2000\ldots$، رقم‌های ۰ و ۹ را هرگز برای $x$ انتخاب نمی‌کنیم. حالا $x$ با هر $f(n)$ دست‌کم در رقم $n$ام فرق دارد، پس $x\neq f(n)$ برای هر $n$. یعنی $f$ پوشا نیست؛ تناقض.
\end{proof}

\section{زبان‌هایی که تشخیص‌پذیر نیستند}

نتیجه قضیه \ref{thm:reals} را این‌طور به کار می‌بریم: نشان می‌دهیم ماشین‌های تورینگ شمارا هستند، اما زبان‌ها ناشمارا هستند. چون هر ماشین فقط یک زبان را تشخیص می‌دهد، زبان‌هایی می‌مانند که هیچ ماشینی تشخیصشان نمی‌دهد.

\begin{corollary}\label{cor:non-rec}
زبان‌هایی وجود دارند که تورینگ-تشخیص‌پذیر نیستند.
\end{corollary}
\begin{proof}
\textbf{مجموعه همه رشته‌ها شماراست.} برای هر الفبای $\Sigma$، رشته‌های $\Sigma^*$ را به ترتیب رشته‌ای فهرست می‌کنیم: اول رشته به طول صفر، بعد رشته‌های به طول یک، بعد طول دو و \dots\ هر رشته جای مشخصی در این فهرست دارد.

\textbf{مجموعه همه ماشین‌های تورینگ شماراست.} هر ماشین تورینگ $M$ یک کدگذاری $\enc{M}$ دارد که یک رشته است. اگر از فهرست همه رشته‌ها، رشته‌هایی را که کدگذاری معتبر یک ماشین نیستند حذف کنیم، فهرستی از همه ماشین‌ها به دست می‌آید.

\textbf{مجموعه همه زبان‌ها ناشماراست.} ابتدا نشان می‌دهیم مجموعه $\mathcal{B}$ همه \textbf{دنباله‌های دودویی نامتناهی} ناشماراست. استدلال دقیقاً مثل قضیه \ref{thm:reals} است: اگر $\mathcal{B}$ شمارا بود، دنباله‌ها را فهرست می‌کردیم و دنباله جدیدی می‌ساختیم که بیت $i$امش با بیت $i$ام دنباله $i$ام فرق دارد:
\begin{ltrfig}
\begin{tabular}{l}
$\underline{0}00000\ldots$\\
$0\underline{1}0100\ldots$\\
$10\underline{0}100\ldots$\\
$111\underline{0}00\ldots$\\
$\ \vdots$
\end{tabular}
\qquad$\Longrightarrow$\qquad $1011\ldots$
\end{ltrfig}
این دنباله در فهرست نیست؛ تناقض.

حالا تناظری بین مجموعه همه زبان‌ها $\mathcal{L}$ و $\mathcal{B}$ می‌سازیم. اگر $\Sigma^*=\{s_1,s_2,s_3,\dots\}$ (به ترتیب رشته‌ای)، هر زبان $A$ یک \textbf{دنباله مشخصه} $\chi_A$ دارد که بیت $i$امش ۱ است اگر $s_i\in A$ و ۰ است اگر $s_i\notin A$:
\begin{ltrfig}
\begin{tabular}{rcccccccccl}
$\Sigma^*=\{$ & $\eps,$ & $0,$ & $1,$ & $00,$ & $01,$ & $10,$ & $11,$ & $000,$ & $001,$ & $\cdots\ \}$\\
$A=\{$ & & $0,$ & & $00,$ & $01,$ & & & $000,$ & $001,$ & $\cdots\ \}$\\
$\chi_A=$ & $0$ & $1$ & $0$ & $1$ & $1$ & $0$ & $0$ & $1$ & $1$ & $\cdots$
\end{tabular}
\end{ltrfig}
تابع $f:\mathcal{L}\to\mathcal{B}$ با $f(A)=\chi_A$ یک‌به‌یک و پوشاست (هر زبان یعنی یک انتخاب دودویی برای هر رشته). چون $\mathcal{B}$ ناشماراست، $\mathcal{L}$ هم ناشماراست.

پس تعداد زبان‌ها بیشتر از تعداد ماشین‌هاست و زبانی وجود دارد که هیچ ماشین تورینگی آن را تشخیص نمی‌دهد.
\end{proof}

\section{اثبات تصمیم‌ناپذیری مسئله پذیرش}

\begin{proof}[قضیه \ref{thm:atm}]
با برهان خلف فرض کنید $H$ یک تصمیم‌گیرنده برای $\ATM$ باشد:
\[H(\enc{M,w})=\begin{cases}\text{accept} & \text{\rl{اگر $M$ رشته $w$ را بپذیرد}}\\ \text{reject} & \text{\rl{اگر $M$ رشته $w$ را نپذیرد.}}\end{cases}\]
با استفاده از $H$ ماشین جدید $D$ را می‌سازیم. $D$ از $H$ می‌پرسد «آیا $M$ وقتی کدگذاری خودش را به عنوان ورودی بگیرد، آن را می‌پذیرد؟» و دقیقاً \textbf{برعکس} عمل می‌کند:
\begin{algo}{$D$}{روی ورودی $\enc{M}$ که $M$ یک ماشین تورینگ است}
\begin{enumerate}
\item $H$ را روی ورودی $\enc{M,\enc{M}}$ اجرا کن.
\item برعکس خروجی $H$ عمل کن: اگر $H$ پذیرفت \textbf{رد کن} و اگر $H$ رد کرد \textbf{بپذیر}.
\end{enumerate}
\end{algo}
(اجرای یک برنامه روی کد خودش عجیب نیست؛ مثلاً یک کامپایلر را می‌توان با خودش کامپایل کرد.) پس
\[D(\enc{M})=\begin{cases}\text{accept} & \text{\rl{اگر $M$ رشته $\enc{M}$ را نپذیرد}}\\ \text{reject} & \text{\rl{اگر $M$ رشته $\enc{M}$ را بپذیرد.}}\end{cases}\]
حالا $D$ را روی کدگذاری خودش، $\enc{D}$، اجرا کنید:
\[D(\enc{D})=\begin{cases}\text{accept} & \text{\rl{اگر $D$ رشته $\enc{D}$ را نپذیرد}}\\ \text{reject} & \text{\rl{اگر $D$ رشته $\enc{D}$ را بپذیرد.}}\end{cases}\]
هر کاری $D$ انجام دهد، مجبور است برعکس آن را انجام دهد؛ تناقض. پس نه $D$ و نه $H$ نمی‌توانند وجود داشته باشند.
\end{proof}

\textbf{خلاصه اثبات:}
\begin{itemize}
\item $H$ رشته $\enc{M,w}$ را دقیقاً وقتی می‌پذیرد که $M$ رشته $w$ را بپذیرد.
\item $D$ رشته $\enc{M}$ را دقیقاً وقتی رد می‌کند که $M$ رشته $\enc{M}$ را بپذیرد.
\item $D$ رشته $\enc{D}$ را دقیقاً وقتی رد می‌کند که $D$ رشته $\enc{D}$ را بپذیرد. تناقض!
\end{itemize}

\begin{note}[قطری‌سازی کجاست؟]
جدولی بکشید که سطرهایش ماشین‌های $M_1,M_2,\dots$ و ستون‌هایش کدگذاری‌های $\enc{M_1},\enc{M_2},\dots$ هستند. خانه $(i,j)$ خروجی $H$ روی $\enc{M_i,\enc{M_j}}$ است. $D$ هم یک ماشین است، پس سطری در این جدول دارد؛ اما طبق تعریف، خروجی $D$ روی $\enc{M_i}$ برعکس \textbf{قطر} جدول است. وقتی به خانه $(D,\enc{D})$ می‌رسیم، این خانه باید برعکس خودش باشد:
\begin{ltrfig}
\small
\begin{tabular}{c|cccccc}
 & $\langle M_1\rangle$ & $\langle M_2\rangle$ & $\langle M_3\rangle$ & $\langle M_4\rangle$ & $\cdots$ & $\langle D\rangle$\\ \hline
$M_1$ & \underline{acc} & rej & acc & rej & & acc\\
$M_2$ & acc & \underline{acc} & acc & acc & & acc\\
$M_3$ & rej & rej & \underline{rej} & rej & & rej\\
$M_4$ & acc & acc & rej & \underline{rej} & & acc\\
$\vdots$ & & & & & $\ddots$ & \\
$D$ & rej & rej & acc & acc & & \textbf{?}
\end{tabular}
\end{ltrfig}
\end{note}

\section{یک زبان تشخیص‌ناپذیر}

$\ATM$ تشخیص‌پذیر است ولی تصمیم‌پذیر نیست. حالا زبانی مشخص پیدا می‌کنیم که حتی تشخیص‌پذیر هم نیست. یادآوری: متمم یک زبان، همه رشته‌هایی است که در آن زبان نیستند.

\begin{definition}[هم-تشخیص‌پذیر]
زبانی \textbf{هم-تورینگ-تشخیص‌پذیر} (\lr{co-Turing-recognizable}) است اگر متمم یک زبان تورینگ-تشخیص‌پذیر باشد.
\end{definition}

\begin{theorem}\label{thm:dec-corec}
یک زبان تصمیم‌پذیر است اگر و تنها اگر هم تورینگ-تشخیص‌پذیر باشد و هم هم-تورینگ-تشخیص‌پذیر.
\end{theorem}
\begin{proof}
\textbf{جهت رفت:} اگر $A$ تصمیم‌پذیر باشد، تشخیص‌پذیر هم هست. همچنین $\overline{A}$ هم تصمیم‌پذیر است (کافی است خروجی تصمیم‌گیرنده $A$ را برعکس کنیم) و در نتیجه تشخیص‌پذیر است. پس $A$ هم-تشخیص‌پذیر است.

\textbf{جهت برگشت:} فرض کنید $A$ و $\overline{A}$ هر دو تشخیص‌پذیر باشند؛ $M_1$ تشخیص‌دهنده $A$ و $M_2$ تشخیص‌دهنده $\overline{A}$ باشد. تصمیم‌گیرنده $M$ برای $A$:
\begin{algo}{$M$}{روی ورودی $w$}
\begin{enumerate}
\item $M_1$ و $M_2$ را روی $w$ \textbf{به‌موازات} اجرا کن.
\item اگر $M_1$ پذیرفت \textbf{بپذیر}؛ اگر $M_2$ پذیرفت \textbf{رد کن}.
\end{enumerate}
\end{algo}
«به‌موازات» یعنی $M$ دو نوار دارد و یکی در میان یک گام از $M_1$ و یک گام از $M_2$ را اجرا می‌کند تا یکی از آن‌ها بپذیرد. (اجرای $M_1$ تا پایان و بعد $M_2$ کار نمی‌کند، چون ممکن است $M_1$ حلقه بزند.)

هر رشته $w$ یا در $A$ است یا در $\overline{A}$. پس دقیقاً یکی از $M_1$ و $M_2$ بالاخره $w$ را می‌پذیرد و $M$ همیشه متوقف می‌شود. اگر $w\in A$، $M_1$ می‌پذیرد و $M$ هم می‌پذیرد؛ اگر $w\notin A$، $M_2$ می‌پذیرد و $M$ رد می‌کند. پس $M$ تصمیم‌گیرنده $A$ است.
\end{proof}

\begin{corollary}
$\overline{\ATM}$ تورینگ-تشخیص‌پذیر نیست.
\end{corollary}
\begin{proof}
می‌دانیم $\ATM$ تشخیص‌پذیر است. اگر $\overline{\ATM}$ هم تشخیص‌پذیر بود، طبق قضیه \ref{thm:dec-corec} $\ATM$ تصمیم‌پذیر می‌شد. اما طبق قضیه \ref{thm:atm} این‌طور نیست. پس $\overline{\ATM}$ تشخیص‌پذیر نیست.
\end{proof}

\begin{summary}
\begin{itemize}
\item $\ATM$ تشخیص‌پذیر است (با ماشین جهانی $U$) اما تصمیم‌پذیر نیست (قطری‌سازی با ماشین $D$).
\item $\mathbb{N}$، اعداد زوج و اعداد گویا شمارا هستند؛ $\mathcal{R}$ و مجموعه دنباله‌های دودویی نامتناهی ناشمارا هستند.
\item ماشین‌ها شمارا و زبان‌ها ناشمارا هستند، پس زبان‌های تشخیص‌ناپذیر وجود دارند.
\item تصمیم‌پذیر = تشخیص‌پذیر و هم-تشخیص‌پذیر. پس $\overline{\ATM}$ تشخیص‌پذیر نیست.
\end{itemize}
\end{summary}
